\section{Selection by the residual entrance phase}
\label{sec:contact-selection}

Section~\ref{sec:graph} derives the entrance phase and its actual
same-index accuracy; Section~\ref{sec:physical-paths} proves convergence
and realization of the complete ordered paths. We now compare the
surviving entrance difference with the actual fixed-angle stopping
window. This gives one selection law, followed by its belt and profile
applications.

Fix $\lambda$, a sufficiently small $0<s\le s_0$, and put
\[
 A_0=9/32,\qquad \rho_c=3/4,\qquad
 S_s=\mathcal S_\lambda(s)>0,\qquad d_*=F_\lambda^{-1}(s).
\]
The entrance constant $\chi_\lambda(\mathfrak a)$ is defined in
\eqref{sel:chi}; it is continuous and bounded on the prescribed compact
auxiliary classes.

\subsection{Selection by the entrance constant}

\begin{proposition}[Selection by the entrance constant]
\label{sel:threshold-selection}
\label{sel:auxiliary-chi}
Under the standing hypotheses the following assertions hold.
\begin{enumerate}[label=\textup{(\roman*)},leftmargin=*]
\item Let $\mathfrak K$ be any nonempty compact auxiliary class within
the prescribed classes. For any $\chi_*\in\R$, define
\begin{equation}
 \eps_m(\chi_*)
 =\exp\{-A_0^{-1}\rho_c^{-m}S_s+A_0^{-1}\chi_*\},
 \qquad j=m+3.
 \label{sel:threshold-sequence}
\end{equation}
This sequence strictly decreases to zero. For every fixed
$\mathfrak a\in\mathfrak K$ with
$\chi_\lambda(\mathfrak a)\ne\chi_*$, its eventual touchdown
count and its terminal-pair and ordered-path limits are
\[
\begin{array}{c|c|c|c}
 &N_s&\text{terminal pair}&\text{ordered path}\\ \hline
 \chi_\lambda(\mathfrak a)<\chi_* &j&(d_*,3)
   &\cW_{\lambda,\nu,s}(s,0)\\
 \chi_\lambda(\mathfrak a)>\chi_* &j+1&(s,2)
   &\cW_{\lambda,\nu,s}(s,2).
\end{array}
\]
For every fixed $\Delta_\chi>0$, the index conclusions hold
simultaneously for all sufficiently large $m$, and the parameter,
ordered-path and Hausdorff limits are uniform, on each of the
compact subclasses
\[
 \{\mathfrak a\in\mathfrak K:
           \chi_\lambda(\mathfrak a)\le\chi_*-\Delta_\chi\},
 \qquad
 \{\mathfrak a\in\mathfrak K:
           \chi_\lambda(\mathfrak a)\ge\chi_*+\Delta_\chi\}.
\]
Each incoming datum is compared with its own arc $I_{\lambda,\nu}$.
No conclusion at equality is asserted.
\item For any two fixed auxiliary labels
$\mathfrak a_1,\mathfrak a_2$ in the prescribed classes with
$\chi_\lambda(\mathfrak a_1)\ne\chi_\lambda(\mathfrak a_2)$,
there is a positive compliance threshold, depending on this fixed
pair, below which their two first section times cannot both occur
after touchdown and at or before its following release.
In particular there is no common compliance sequence selecting
contact interiors for both labels, either in terminal parameters or
in ordered paths, even at different phases.

\item Suppose that, for a fixed label $\mathfrak a_0$ and
$0<\alpha_0<2$, a sequence $\eps_m\downarrow0$ satisfies
\[
 f_{\eps_m}^{\mathfrak a_0}
 \longrightarrow\cW_{\lambda,\nu_0,s}(s,\alpha_0)
 \qquad\hbox{in }\dF.
\]
Along this same given sequence, every fixed label $\mathfrak a$ with
unequal entrance constant has the following limits:
\[
\begin{array}{c|c}
 &\text{terminal pair; ordered path}\\ \hline
 \chi_\lambda(\mathfrak a)<\chi_\lambda(\mathfrak a_0)
   &(d_*,3);\quad\cW_{\lambda,\nu,s}(s,0)\\
 \chi_\lambda(\mathfrak a)>\chi_\lambda(\mathfrak a_0)
   &(s,2);\quad\cW_{\lambda,\nu,s}(s,2).
\end{array}
\]
This assertion uses the given reference sequence; it makes no
claim about its touchdown count. The corresponding Hausdorff image
convergence holds, and each path retains its own incoming arc
$I_{\lambda,\nu}$.
Equality of the entrance constants is necessary for common interior
contact selection; its sufficiency is not determined by these estimates.
\end{enumerate}
\end{proposition}

Taking $\chi_*>\max_{\mathfrak K}\chi_\lambda$ or
$\chi_*<\min_{\mathfrak K}\chi_\lambda$ gives explicit common
sequences for the admission or release endpoint, respectively.
The threshold need not be attained by a label. These sequences use
the exact phase at $s$, as the free-target sequence in
Theorem~\ref{thm:common-free} uses the exact phase at its target.
The scale comparison is
\[
 \eps\ll1/\log(1/\eps)\ll1:
\]
the retained entrance separation vanishes, but exceeds the actual
contact stopping window. Figure~\ref{fig:selection-mechanism} shows
the retained phase gap and the actual stopping window.
% MIXED: schematic base positions with exact threshold and window relations.
% The horizontal coordinate is the actual base at one common event index.
\begin{figure}[tbp]
\centering
\begin{tikzpicture}[
 x=1cm,y=1cm,font=\small,
 every node/.style={inner sep=1.5pt},
 arrow/.style={-{Latex[length=1.6mm,width=1.1mm]},line width=.5pt},
 brace/.style={decorate,decoration={brace,mirror,amplitude=3pt},
                line width=.45pt}
]
 \node[anchor=west,font=\small\bfseries] at (0,4.80)
   {Residual phase along the threshold sequence, $j=m+3$};
 \node at (7.15,4.18)
   {$\displaystyle \mathcal S_\lambda(\delta_j)-\mathcal S_\lambda(s)
    =\rho_c^m\{\chi-\chi_*+r_{\mathfrak a}(\varepsilon_m)\}
       +O(\varepsilon_m^{3/32})$};
 \node at (7.15,3.56)
   {$\displaystyle \rho_c^m\sim
       \frac{\mathcal S_\lambda(s)}{A_0L},\qquad
       A_0=\frac9{32},\quad\rho_c=\frac34,\quad
       L=\log(1/\varepsilon_m)$};
 \node[align=center] at (2.30,2.80)
   {$\chi>\chi_*$\\release: $N_s=j+1$};
 \node[align=center] at (12.00,2.80)
   {$\chi<\chi_*$\\admission: $N_s=j$};
 \fill[black!13] (6.60,1.70) rectangle (7.35,2.12);
 \draw[arrow] (.40,1.91) -- (14.10,1.91)
   node[right] {$\delta_j$};
 \draw[densely dotted,line width=.5pt] (6.60,1.58) -- (6.60,2.22);
 \draw[line width=.6pt] (7.35,1.58) -- (7.35,2.22);
 \node[above right] at (7.35,2.13) {$s$};
 \fill (2.30,1.91) circle[radius=1.9pt];
 \fill (12.00,1.91) circle[radius=1.9pt];
 \draw[brace] (2.30,1.44) -- (7.35,1.44);
 \draw[brace] (7.35,1.44) -- (12.00,1.44);
 \node at (4.36,1.07) {$\Theta_s(L^{-1})$};
 \node at (9.765,1.07) {$\Theta_s(L^{-1})$};
 \node[align=center] at (6.975,.79)
   {$O_s(\varepsilon_m)$\\contact window};
 \node[align=center] at (2.30,.16)
   {$(s,2)$\\$\mathcal W_{\lambda,\nu,s}(s,2)$};
 \node[align=center] at (12.00,.16)
   {$(d_*,3)$\\$\mathcal W_{\lambda,\nu,s}(s,0)$};
 \node at (7.15,-.65)
   {$\varepsilon_m\ll L^{-1}\ll1,
       \qquad O(\varepsilon_m^{3/32})=o(L^{-1})$};
 \node[align=center] at (7.15,-1.22)
   {At $\chi=\chi_*$ the leading comparison does not determine the stopping phase.};
\end{tikzpicture}
\caption{Residual phase versus contact resolution, with $\lambda,s$
fixed, and $\chi=\chi_\lambda(\mathfrak a)$.
The exact phase decreases with the base, giving opposite signs
for the phase and base gaps. On compact label subclasses separated
from $\chi_*$ by a fixed margin, the base gaps have order $L^{-1}$,
larger than both the comparison error and the necessary contact-stop
window $0<s-\delta_j\le C_s\varepsilon_m$.
The displayed pairs are the actual terminal-pair limits.
The admission-side touchdown is read on the original trajectory
continued past observation. Positions and widths are schematic;
no label attaining $\chi_*$ is assumed.}
\label{fig:selection-mechanism}
\end{figure}


The following two lemmas connect this phase comparison to actual
stopping events and complete path limits.

\begin{lemma}[Identification of an interior contact target]
\label{sel:path-identification}
For fixed labels, if $f_{\eps_m}\to\cW(s,\alpha)$ in $\dF$, with
$0<\alpha<2$, then
$(d_{\eps_m},\alpha_{\eps_m})\to(s,\alpha)$.
\end{lemma}
\begin{proof}
Every subsequence has a further convergent terminal-pair subsequence,
by Theorem~\ref{thm:whole}. If its limit is $(d,\beta)\in\mathcal A_s$,
that theorem gives the path limit $\cW(d,\beta)$.
Ordered convergence preserves endpoints, hence
$\cC_d(\beta)=\cC_s(\alpha)$.
An interior free endpoint has positive normal gap; on contact
$y=z=0$. The two free endpoints are contact admission and release.
The contact endpoint map is injective on $[0,2]$ by
\eqref{gr:contact-injectivity}, including the arrest point $\beta=1$.
An interior contact point is neither admission nor release, so neither
free endpoint nor the repeated admission representation can give this
endpoint. The only possible pair is $(s,\alpha)$, proving convergence
of the full parameter sequence.
\end{proof}

\begin{lemma}[Endpoint test at a fixed touchdown index]
\label{sel:endpoint-test}
Suppose that along a compliance sequence the actual $j$th touchdown
is defined within the common input domain, $j\to\infty$, and
$\delta_j\to s$. If $(\delta_j-s)/\eps\to+\infty$, then
$N_s=j$, the terminal pair tends to $(d_*,3)$ and the ordered path
tends to $\cW_{\lambda,\nu,s}(s,0)$. If
$(\delta_j-s)/\eps\to-\infty$, then $N_s=j+1$, the pair tends
to $(s,2)$ and the path tends to $\cW_{\lambda,\nu,s}(s,2)$.
The index assertions are eventual. All limits, including the
Hausdorff image limits, are uniform on a compact auxiliary class
whenever the hypotheses are uniform there.
\end{lemma}
\begin{proof}
The growth bounds give a positive lower bound for the predecessor
base. Its complete state tends to the graph by
Lemma~\ref{path:stopped-cascade}; the actual return and strict
increase of $F_\lambda$ then give
$\delta_{j-1}\to F_\lambda^{-1}(s)=d_*<s$.
For a positive gap at index $j$, the predecessor contact stays below
the section: its base remains a fixed distance below $s$ and its
contact drift is $O_s(\eps)$. Since $\delta_j>s$, the first
crossing is in its following free flight, before touchdown $j$.
Thus $N_s=j$ and $d_\eps\to d_*$. The unique solution of the
reference section equation at $d_*$ is the free endpoint $3$.

For a negative gap, contact $j$ changes the base coordinate by
$O_s(\eps)$, less than $s-\delta_j$, so it cannot reach the
section. Its input is below the inner cutoffs and has all the
admission margins. Apply the actual return estimate at this input;
its next output tends to $F_\lambda(s)>s$. The first crossing is
in that free flight, with $N_s=j+1$, $d_\eps\to s$ and
$\alpha_\eps>2$. The reference section equation at $s$ has the
unique solution $2$ on $[2,3]$. This use of the next return requires
no iteration at its output.

Strict decrease \eqref{path:strict-p} and the original-prefix
control in Lemma~\ref{path:first-section} exclude earlier crossings.
Theorem~\ref{thm:whole} gives the path limits, using
$\cW(d_*,3)=\cW(s,0)$ at admission. For uniform limits, all base
bounds, contact drifts and complete-state budgets above are uniform.
A varying-label subsequence can have only the stated unique free
root. Match any fixed finite number of last returns uniformly and
put the entire preceding cascade in its common small ball, as in
\eqref{path:whole-error}. The incoming prefix is compared with its
own $I_{\lambda,\nu}$. This gives uniform ordered convergence, and
\eqref{eq:ps5-hausdorff} gives the uniform image convergence.
\end{proof}

\begin{proof}[Proof of Proposition~\ref{sel:threshold-selection}]
Write $L=\log(1/\eps)$. For \textup{(i)}, the exponent in
\eqref{sel:threshold-sequence} strictly decreases to $-\infty$, and
$A_0L=\rho_c^{-m}S_s-\chi_*$. The graph chain started at the actual
$\delta_3^{\mathfrak a}$ has phase at $j=m+3$ equal to
\[
 S_s+\rho_c^m\{\chi_\lambda(\mathfrak a)-\chi_*
                         +r_{\mathfrak a}(\eps_m)\}.
\]
The bounded entrance constants and uniform remainder in
\eqref{sel:uniform-entry-phase} put these graph bases uniformly
near $s$; all bases through $j$ are eventually below $3s/2$.
Lemma~\ref{sel:sharp-phase} therefore continues the actual events
through this same index for every label and gives
\begin{equation}
 \mathcal S_\lambda(\delta_j^{\mathfrak a})-S_s
 =\rho_c^m\{\chi_\lambda(\mathfrak a)-\chi_*
                        +r_{\mathfrak a}(\eps_m)\}
                       +O(\eps_m^{3/32}).
 \label{sel:threshold-phase}
\end{equation}
Here $\rho_c^m\sim S_s/(A_0L)$ and the last error is
$o(\rho_c^m)$. On either compact subclass separated from $\chi_*$
by $\Delta_\chi>0$, the phase difference has the sign of
$\chi_\lambda(\mathfrak a)-\chi_*$ and absolute value between
$c\Delta_\chi\rho_c^m$ and $C\rho_c^m$, uniformly.
Strict decrease and the two-sided Lipschitz bound
\eqref{gr:phase-bilip} give the opposite signed base gap of the
same order. Since $\eps_m=o(\rho_c^m)$,
Lemma~\ref{sel:endpoint-test} proves all index, terminal-pair and
path assertions, including uniformity and image convergence.
Every fixed unequal label belongs to such a separated subclass.
At equality the sign is not determined.

For the event exclusion in \textup{(ii)}, suppose instead that
two fixed unequal labels both stop during contact along a common
sequence tending to zero, at touchdown indices $j_1,j_2$.
The contact window \eqref{eq:contact-stop-location} and the sharp phase
give, for $i=1,2$,
\[
 \rho_c^{j_i-3}
       \{A_0L+\chi_\lambda(\mathfrak a_i)+o(1)\}
                =S_s+O_s(\eps^{3/32}).
\]
Thus $\rho_c^{j_i-3}A_0L\to S_s$. Division gives
$\rho_c^{j_1-j_2}\to1$; because the index difference is an integer,
eventually $j_1=j_2=j$. Subtracting the two equations gives
\[
 \rho_c^{j-3}
 \{\chi_\lambda(\mathfrak a_1)-\chi_\lambda(\mathfrak a_2)+o(1)\}
                   =O_s(\eps^{3/32}).
\]
Multiplication by $L$ contradicts
$\rho_c^{j-3}L\to S_s/A_0>0$ and the unequal constants.
Consequently the exclusion holds below a positive threshold for
this fixed pair. Interior terminal-pair limits eventually give
actual contact stops; Lemma~\ref{sel:path-identification} gives
the same implication for interior ordered-path limits.

For \textup{(iii)}, the identification lemma
puts the $\mathfrak a_0$ stops in contact. Set
$j=N_s^{\mathfrak a_0}-1$. Then
\[
 \delta_j^{\mathfrak a_0}=s+O_s(\eps),\qquad
 \rho_c^{j-3}\sim\frac{S_s}{A_0L}.
\]
For any other fixed $\mathfrak a$, the graph chain started at its
own actual $\delta_3^{\mathfrak a}$ has phase at this index tending
to $S_s$. Its bases through $j$ are eventually below $3s/2$.
The continuation in Lemma~\ref{sel:sharp-phase} therefore compares
the actual events through the same index, even after that
trajectory's observation time. Subtracting the reference relation
gives
\begin{equation}
 \mathcal S_\lambda(\delta_j^{\mathfrak a})-S_s
 =\frac{S_s}{A_0L}
   \{\chi_\lambda(\mathfrak a)
       -\chi_\lambda(\mathfrak a_0)+o(1)\}
       +O_{\mathfrak a,\mathfrak a_0,s}(\eps^{3/32}).
 \label{sel:separation}
\end{equation}
The signed base gap is opposite to the fixed entrance difference
and has size $\Theta(L^{-1})$. Lemma~\ref{sel:endpoint-test} gives
the terminal-pair and ordered-path limits in \textup{(iii)}, with
the stated incoming arcs and image convergence. It is applied at the
index fixed by the reference label, not at the prescribed index of
\textup{(i)}, so no touchdown-count formula for the given sequence is
asserted.
\end{proof}

\begin{proof}[Proof of Theorem~\ref{thm:contact-selection}]
Proposition~\ref{loc:entrance} constructs the entrance constant from
the physical preparation and proves its continuity and boundedness.
The same-index comparison and actual stopping-window estimates lead
to Proposition~\ref{sel:threshold-selection}. Its part
\textup{(ii)} gives assertion \textup{(i)} of the theorem, and
part \textup{(iii)} gives assertion \textup{(ii)}. The common compact classes
and each label's incoming arc are exactly those of
Section~\ref{sec:model}.
\end{proof}

\begin{remark}[Phase and stopping resolution]
\label{sel:criterion}
The integer-index argument uses the following general hypotheses.
Let $L=\log(1/\eps)$, $A>0$, $S_*>0$ and $0<\rho<1$.
Suppose that at the relevant integer event indices
$n\in\mathbb N_0$, for each fixed label $\eta$,
\[
 P_n(\eps,\eta)
   =\rho^n\{AL+\chi(\eta)+r_\eta(\eps)\}+R_n(\eps,\eta),
 \qquad r_\eta\to0,\quad R_n=o(L^{-1}),
\]
where the estimate for $R_n$ is uniform over all relevant event
indices for that fixed label. If the event requires
$|P_n-S_*|\le w_\eta(\eps)=o(L^{-1})$, two fixed labels with
unequal $\chi$ cannot both meet it for arbitrarily small common
$\eps$. The ratio and subtraction argument in the proof of
Proposition~\ref{sel:threshold-selection}\textup{(ii)} applies with
$n_i$ in place of $j_i-3$.
The actual preparation, same-index accuracy and stopping window
needed here were established in Sections~\ref{sec:local-results}
and~\ref{sec:graph}.
\end{remark}

\subsection{Belt ordering of the entrance constant}
\label{sec:entrance-phase}

Theorem~\ref{thm:contact-selection} reduces belt dependence to the
order of the prepared coefficient $C_R$. The following strict order
gives Corollary~\ref{cor:belt-contact-selection}.
Throughout this subsection $\lambda,\theta$ are fixed; the mechanical
coefficient is denoted by $d_\lambda$ as in Section~\ref{sec:model}.

\begin{proposition}[Belt order of the entrance constant]
\label{sel:entrance-phase}
Let $C_R(k)$ be the second-touchdown coefficient in
\eqref{loc:preparation-constants}, and let $K_\lambda$ be the
return coefficient in \eqref{loc:constants}. Then $C_R(k)$ strictly decreases on $I$. Hence
$\chi_{\lambda,\theta}(k)$ in \eqref{sel:chi} strictly increases.
\end{proposition}
\begin{proof}
The normalized critical passage determines $q=q_{\rm in}>0$ from
$(\lambda,\theta)$ alone; see Lemma~\ref{lem:critical-pullback}
and Theorem~\ref{thm:detuned-release}. The normal-release constant
$\nu_\#$ is determined by the fixed linear problem
\eqref{con:canonical-release}, independently of the original labels.
Put $H=ak/\gamma$ and $t=T_f(k)$ in
\eqref{loc:preparation-constants}, where $A_{\rm reb}$ is the
amplitude $A_s$ of \eqref{prep:amplitude}.
Implicit differentiation of the flight equation gives
\[
 t'(k)=-\frac{3\ell(a/\gamma)t}{2ct+3\ell H}<0.
\]
Substitution in the amplitude first yields
\[
 A_{\rm reb}^2=\frac{d_\lambda}{2h^2}
 \left[\gamma Ht^2+(2hc/3+ad_\lambda)t^3\right].
\]
Eliminating $H=(6q-ct^2)/(3\ell t)$ and using
$\gamma=h\ell-d_\lambda$, we obtain
\begin{equation}
 A_{\rm reb}^2=\frac{d_\lambda}{2h^2}
 \left[\frac{2\gamma q}{\ell}t+
 \left(\frac{hc}{3}+\frac{d_\lambda c}{3\ell}+ad_\lambda\right)t^3\right].
 \label{sel:positive-amplitude}
\end{equation}
Every coefficient is positive. Hence the amplitude, $R$ and $C_R$
strictly decrease with $k$, which proves the order of $\chi$.

\end{proof}

The profile dependence requires a different argument, because a
slip level can be visited more than once. We next construct a
single-visit perturbation with nonzero release-speed response.


\begin{proof}[Proof of Corollary~\ref{cor:belt-contact-selection}]
For the exclusion, take $\mathfrak a_i=(k_i,\theta,\nu,\mathbf n)$.
Proposition~\ref{sel:entrance-phase} gives unequal entrance constants
for distinct fixed $k_i$; Proposition~\ref{sel:auxiliary-chi}\textup{(ii)}
therefore gives the event, terminal-pair and ordered-path exclusions.
For a given reference sequence, apply Proposition~\ref{sel:auxiliary-chi}\textup{(iii)} with
$\mathfrak a_0=(k_0,\theta,\nu,\mathbf n)$ and
$\mathfrak a=(k,\theta,\nu,\mathbf n)$. The strict increase of
$\chi_{\lambda,\theta}(k)$ gives admission for $k<k_0$ and release
for $k>k_0$, with the same incoming arc and the corresponding
Hausdorff image convergence.
\end{proof}

\subsection{Friction-profile response}
\label{sec:profile-response}

The belt order does not by itself establish sensitivity to the friction
profile. That question is decided by the normalized critical passage.
A slip level can be visited several times, so a change of the profile
does not provide an independently chosen forcing at each visit.
We use an incoming slip interval visited only once before first release
and prove that its release-speed response cannot vanish in every
admissible direction.

\begin{proposition}[Profile response and entrance order]
\label{sel:profile-response}
Fix a mechanical parameter satisfying the critical-rate and sign
conditions of Section~\ref{sec:model}, and an admissible $C^4$
profile $\Phi^{\rm b}$. Write $q_{\rm in,\lambda}[\Phi]=q_{\rm in}$ for the
normalized first-release velocity in Lemma~\ref{lem:critical-pullback}.
There is $H>Q_\Phi$
such that, for every nonempty open interval $J\subset(H,\infty)$,
there are a nonzero $\psi\in C_c^\infty(J)$ and
$\vartheta_{\max}>0$ for which
\[
 \{\Phi^{\rm b}+\vartheta\psi:
                        |\vartheta|\le\vartheta_{\max}\}
\]
is an admissible compact profile class and
$D_\Phi q_{\rm in,\lambda}[\Phi^{\rm b}](\psi)\ne0$.
Constructing the curve at $\lambda_c=(7/3,0)$, we may orient it
and choose a sufficiently small relative closed mechanical box
$\Lambda_{\rm fr}$ about $\lambda_c$ and a closed parameter interval so that
\begin{equation}
 \partial_\vartheta q_{\rm in,\lambda}[\Phi_\vartheta]\ge c_{\rm fr}>0
 \quad\text{on }\Lambda_{\rm fr}\times
                   [-\vartheta_{\max},\vartheta_{\max}],
 \qquad \Phi_\vartheta=\Phi^{\rm b}+\vartheta\psi.
 \label{sel:profile-local-order}
\end{equation}
For each fixed $\lambda\in\Lambda_{\rm fr}$ and $k\in I$,
$\chi_{\lambda,\vartheta}(k)$ strictly decreases with $\vartheta$.
\end{proposition}
\begin{proof}
On the baseline critical contact write $L=Z-Y$, let $\tau_r$
be its first release, and abbreviate $C=C_G$, $m=\mu A_G$ and
$f=1-\Phi^{\rm b}$. The normalization gives $U\to+\infty$ and
$U'\to-\gamma<0$ as $\tau\to-\infty$. Choose $\tau_0<\tau_r$
with $U'<-\gamma/2$ on $(-\infty,\tau_0]$, and take
\[
 H>\max\{Q_\Phi,\max_{[\tau_0,\tau_r]}U\}.
\]
Every level $u>H$ is visited exactly once before first release;
denote its inverse time by $\tau(u)$.

Let $F_{\Phi^{\rm b}}$ be the contact extension of
\eqref{ent:critical-field}, obtained by replacing $L_+$ by $L$.
In state order $(Y,V,Z,U)$ put
$A(\tau)=D_WF_{\Phi^{\rm b}}(\overline W(\tau))$ and solve
\[
 \zeta'=-A^T\zeta,\qquad
 \zeta(\tau_r)=e_V=(0,1,0,0)^T,\qquad Q=C\zeta_V+m\zeta_U.
\]
For an admissible compactly supported variation,
Lemma~\ref{lem:profile-variation} gives the exact response
\begin{equation}
 D_\Phi q_{\rm in,\lambda}[\Phi^{\rm b}](\psi)
   =-\int_{-\infty}^{\tau_r}LQ\psi(U)\,d\tau.
 \label{sel:profile-adjoint}
\end{equation}
A controllability calculation prevents $Q$ from vanishing on an
interval. Indeed,
\begin{equation}
 \begin{gathered}
 A=A_0+b_0r^T,\qquad
 A_0=\begin{pmatrix}
 0&1&0&0\\0&0&0&0\\1&0&-2&0\\h&0&-h&0
 \end{pmatrix},\\
 b_0=(0,C,0,m)^T,\qquad r=(-f,0,f,f'L)^T,\qquad Q=b_0^T\zeta,
 \\
 \det[b_0,A_0b_0,A_0^2b_0,A_0^3b_0]=-hC^4\ne0.
 \end{gathered}
 \label{sel:profile-controllability}
\end{equation}
If $Q$ vanished on an open contact interval, then
$\zeta'=-A_0^T\zeta$ there. Successive differentiation of
$b_0^T\zeta=0$ would give
$(A_0^jb_0)^T\zeta=0$ for $j=0,1,2,3$.
The determinant in \eqref{sel:profile-controllability} would force
$\zeta=0$, contradicting its nonzero terminal value by linear
uniqueness. Thus $Q$ is not identically zero on any open contact
interval; no analyticity is needed.

Take a nonempty open $J\subset(H,\infty)$. Strict monotonicity,
the mean-value theorem and continuity of the derivative give an open
subinterval where $(\Phi^{\rm b})'>0$. Its time preimage contains
a point with $Q\ne0$. We may therefore choose a closed $K\Subset J$
on which $(\Phi^{\rm b})'$ has a positive lower bound and
$Q(\tau(u))$ has a fixed nonzero sign. Changing variables in
\eqref{sel:profile-adjoint} gives
\[
 D_\Phi q_{\rm in,\lambda}[\Phi^{\rm b}](\psi)
   =-\int_K\frac{L(\tau(u))Q(\tau(u))}{|U'(\tau(u))|}
                   \psi(u)\,du.
\]
A nonzero nonnegative smooth bump supported inside $K$ has nonzero
response. Reverse its sign if needed to make the response positive.

For sufficiently small $|\vartheta|\le\vartheta_{\max}$, the
profile derivative stays positive on $K$, while the profile is
unchanged elsewhere. Its zero, range and saturation limits are
preserved. The core strip and (F3) are unchanged because
$K\subset(Q_\Phi,\infty)$; (F4) holds with a common enlarged
constant since the perturbation has fixed compact support. The
closed parameter curve is continuous in the weighted $C^4$ metric,
and hence compact.

For the common mechanical range choose this curve first at
$\lambda_c$. The model restrictions provide a common mechanical
box for the fixed compact class. Uniform incoming tails give a common
finite time before the support is reached. At each fixed $\lambda$,
all profiles have exactly their baseline's state at that time, so
their shape derivative starts at zero. The baseline state is
continuous in $\lambda$ by Lemma~\ref{lem:critical-pullback}.
Continuous dependence of the finite-time state and linear variation,
together with the transverse first-release event, makes
$\partial_\vartheta q_{\rm in,\lambda}[\Phi_\vartheta]$ jointly continuous
in $(\lambda,\vartheta)$. Its positive value at $(\lambda_c,0)$
gives \eqref{sel:profile-local-order} after restricting the box and
parameter interval. No derivative with respect to $\lambda$ is needed.

Finally keep $\lambda,k$ fixed, put $H_k=ak/\gamma$, and let
$t=t(q)>0$ solve
\[
 ct^2+3\ell H_kt=6q,\qquad t_q=\frac6{2ct+3\ell H_k}>0.
\]
The first amplitude identity in the proof of
Proposition~\ref{sel:entrance-phase} gives
$A_{\rm reb}^2=d_\lambda B(t)/(2h^2)$, where
\[
 B(t)=\gamma H_kt^2+(2hc/3+ad_\lambda)t^3.
\]
As $R=\nu_\#A_{\rm reb}$ and $C_R=\sqrt{6R/c}$,
\begin{equation}
 \partial_q\chi=-\frac3{16}\frac{B'(t)}{B(t)}t_q<0.
 \label{sel:profile-entrance-order}
\end{equation}
Here $\nu_\#$ and $K_\lambda$ are independent of the original
profile. Combining \eqref{sel:profile-local-order} and
\eqref{sel:profile-entrance-order} proves the strict entrance order.
\end{proof}

\begin{corollary}[Friction-profile sensitivity at a fixed belt coefficient]
\label{cor:friction-profile-selection}
Let $\Phi^{\rm b}$ be an admissible $C^4$ baseline profile, and fix
compact belt and physical-preparation classes as in
Section~\ref{sec:model}. There exist a nonzero
$\psi\in C_c^\infty((0,\infty))$, $\vartheta_{\max}>0$, a
sufficiently small relative closed mechanical box $\Lambda_{\rm fr}$
about $\lambda_c=(7/3,0)$, and $s_{\rm fr}>0$ such that
\[
 \Phi_\vartheta=\Phi^{\rm b}+\vartheta\psi,\qquad
 \vartheta\in\Theta_{\rm fr}
              =[-\vartheta_{\max},\vartheta_{\max}]
\]
is a compact class satisfying (F1)--(F4). With this curve as the
profile class, the common model restrictions and
Theorem~\ref{thm:whole} apply on $\Lambda_{\rm fr}$ for
$0<s\le s_{\rm fr}$. Orient the curve so that its normalized
first-release speed $q_{\rm in,\lambda}[\Phi_\vartheta]$ increases with
$\vartheta$ for every $\lambda\in\Lambda_{\rm fr}$.

Fix $\lambda\in\Lambda_{\rm fr}$, $k\in I$, $\nu$, $\mathbf n$
and $0<s\le s_{\rm fr}$. In every system retain
$U_b=U_*-k\sqrt\eps$ and vary only the fixed profile parameter.
\begin{enumerate}[label=\textup{(\roman*)},leftmargin=*]
\item For every two distinct fixed
$\vartheta_1,\vartheta_2\in\Theta_{\rm fr}$
there is $\eps_0>0$, depending on this pair and the fixed data,
such that for $0<\eps<\eps_0$ the two actual first section times
cannot both occur after touchdown and at or before its following
release. Consequently no common sequence $\eps_m\downarrow0$
has terminal-pair limits $(s,\alpha_i)$ for both profiles with
$0<\alpha_i<2$. There is likewise no common sequence with ordered
limits $\cW_{\lambda,\nu,s}(s,\alpha_i)$, even if
$\alpha_1\ne\alpha_2$.
\item Fix $\vartheta_0\in\Theta_{\rm fr}$ and $0<\alpha<2$.
Along any sequence $\eps_m\downarrow0$ for which
$f_{\eps_m}^{(\vartheta_0)}\to\cW_{\lambda,\nu,s}(s,\alpha)$
in $\dF$, every other fixed $\vartheta\in\Theta_{\rm fr}$ satisfies
\[
\begin{array}{c|c|c}
 &\text{terminal pair limit}&\text{ordered-path limit}\\ \hline
 \vartheta>\vartheta_0 &(d_*,3)&\cW(d_*,3)=\cW(s,0)\\
 \vartheta<\vartheta_0 &(s,2)&\cW(s,2).
\end{array}
\]
\end{enumerate}
\end{corollary}

This is the locally constructed profile application of
Theorem~\ref{thm:contact-selection}: the curve and its common
mechanical range are constructed together;
no nonconstancy along a prescribed profile family is assumed.
Each profile is fixed before the compliance limit. With
$\mathbf n=0$ the full physical initial state is the same for all
profiles and compliances, while the belt still follows the stated
joint scaling. Theorem~\ref{thm:whole} supplies the reference
sequences in part \textup{(ii)}; the ordered-path limits also give
the corresponding Hausdorff image convergence by
\eqref{eq:ps5-hausdorff}.



\begin{proof}[Proof of Corollary~\ref{cor:friction-profile-selection}]
Construct the curve at $\lambda_c$ by
Proposition~\ref{sel:profile-response}, choose the common model
range for this compact class, and restrict it as in
\eqref{sel:profile-local-order}. Fix the labels in the corollary.
Equation~\eqref{sel:profile-entrance-order} shows that their entrance
constant strictly decreases with $\vartheta$. Apply
Proposition~\ref{sel:auxiliary-chi}\textup{(ii)} with
$\mathfrak a_i=(k,\vartheta_i,\nu,\mathbf n)$.
Part \textup{(ii)} gives both the terminal-pair and ordered-path
exclusions. For the given reference sequence,
Proposition~\ref{sel:auxiliary-chi}\textup{(iii)} gives admission for
$\vartheta>\vartheta_0$ and release for $\vartheta<\vartheta_0$.
Theorem~\ref{thm:whole} supplies actual reference sequences, and
the same comparison gives Hausdorff image convergence.
\end{proof}

\subsection{Return count and contact-root asymptotics}

\begin{corollary}[Double-logarithmic return count]
\label{sel:return-count}
The actual touchdown count satisfies
\begin{equation}
 N_s(\eps)=\frac{\log\log(1/\eps)}{\log(4/3)}+O_s(1).
 \label{sel:count}
\end{equation}
\end{corollary}
\begin{proof}
The last input belongs to $[c_s,s)$, on which the positive phase
has fixed positive upper and lower bounds. Formula
\eqref{sel:sharp-phase-formula} at $n=N_s-1$ therefore gives
$\rho_c^{N_s-4}L\asymp_s1$. Taking logarithms proves the claim.
\end{proof}

For any actual contact-root family $\eps_j$ with the complete label
fixed, write $C_R=C_R(\lambda,k,\theta)$ and let
$j\to\infty$ be the contact number. The same formula and
\eqref{eq:contact-stop-location} give
\begin{align}
 \log(1/\eps_j)
 &=\tfrac{32}{9}\rho_c^{-(j-3)}\mathcal S_\lambda(s)
       -\tfrac{32}{9}\chi_\lambda(\mathfrak a)+o(1),\nonumber\\
 \eps_j
 &=K_\lambda^{32/3}C_R^{-8/3}
 \exp\{-\tfrac{32}{9}\rho_c^{-(j-3)}\mathcal S_\lambda(s)\}
 (1+o(1)).
 \label{sel:contact-roots}
\end{align}
Indeed $\rho_c^{-(j-3)}=O(L_j)$, so division of the shadow error
costs only $O(L_j\eps_j^{3/32})=o(1)$.
Different interior phases share this asymptotic formula; their
individual selection is supplied by the actual scan of
Lemma~\ref{path:continuous-sweep}.


